\documentclass[11pt,a4paper]{article}
\usepackage[T1]{fontenc}
\usepackage{isabelle,isabellesym}
\usepackage{amsmath,amssymb}
\usepackage{pdfsetup}

\urlstyle{rm}
\isabellestyle{it}

\begin{document}

\title{Cross-Domain State Preservation Functor}
\author{Jinwook Kim (Jay)}
\maketitle

\begin{abstract}
We present a generic, reusable framework for verifying that state transitions
are faithfully preserved when synchronized across heterogeneous domains
(blockchains, off-chain ledgers, etc.), and that the synchronization process
remains live in the presence of Byzantine faults.

The framework combines two complementary properties.  For \emph{safety},
the core abstraction is a hierarchy of Isabelle/HOL locales:
\texttt{state\_machine} defines finite-state deterministic transition systems
with terminal states;
\texttt{state\_preservation} captures the naturality (homomorphism) condition
ensuring that a synchronization map commutes with transitions;
\texttt{symmetric\_state\_preservation} extends this to bidirectional
synchronization with inverse roundtrip guarantees; and
\texttt{multi\_domain\_preservation} generalizes to $N$ domains sharing the
same abstract state machine, with a cross-domain consistency theorem.
For \emph{liveness} under Byzantine faults, three additional locales are
introduced:
\texttt{priority\_system} guarantees deterministic resolution of conflicting
requests via a total order on priority keys;
\texttt{deadlock\_free\_locking} guarantees that timeout-based forced release
prevents indefinite locking; and
\texttt{fair\_leader\_system} guarantees starvation freedom under fair leader
scheduling, formulated as an assume-guarantee abstraction of probabilistic
VRF-based leader election.

As a practical application, we instantiate this framework with the
regulatory state model of Oraclizer, a state synchronization oracle for
tokenized assets across blockchain networks and off-chain ledgers.
The model comprises five states (\textsc{Active}, \textsc{Frozen},
\textsc{Seized}, \textsc{Confiscated}, \textsc{Restricted}) and seven
actions.  The safety instantiation proves that a synchronization protocol
with preemptive locking preserves cross-chain regulatory consistency (the
\emph{regulatory homomorphism} theorem), that the global validity invariant
is maintained across synchronization events (\emph{valid state preservation}),
and that the generic multi-domain locale applies parametrically to the
regulatory model.  The liveness instantiation models a Byzantine consensus
system with $f < n/3$ malicious nodes, instantiates the three liveness
locales with regulatory authority levels and a four-component priority key,
and combines safety and liveness in a single
\emph{combined safety liveness} theorem that discharges the honest-node
assumption of the safety proof under Byzantine faults.

The design is inspired by Lochbihler's Merkle Functor (AFP entry
\texttt{ADS\_Functor}), which abstracted authenticated data structures into
composable building blocks.  This entry extends that pattern to a different
axis of abstraction: cross-domain state preservation targeting
synchronization correctness rather than data structure integrity, together
with the liveness guarantees needed for that synchronization to actually
make progress under Byzantine faults.
\end{abstract}

\tableofcontents

\section{Overview}

This AFP entry consists of four theory files organized as two
``generic locale + regulatory instance'' pairs, one for safety and one
for liveness:

\begin{description}
\item[\texttt{State\_Preservation.thy}] defines domain-independent locales
  for state machines, pairwise state preservation (with naturality),
  symmetric bidirectional preservation, and multi-domain preservation.
  All locales are parameterized over abstract types and can be instantiated
  with any domain that has finite states, deterministic transitions, and
  optional terminal states.

\item[\texttt{Regulatory\_Instance.thy}] instantiates the generic safety
  framework with a concrete regulatory state model.  It defines the five
  regulatory states and seven regulatory actions, proves fundamental
  properties (terminal absorption, universal reachability of confiscation,
  transition exclusions with legal justification), models a synchronization
  protocol with preemptive locking, and proves the main safety theorems:
  \emph{regulatory homomorphism} (cross-chain consistency after sync),
  \emph{valid state preservation} (sync maintains the global validity
  invariant), and \emph{multi-domain parametric instantiation}.

\item[\texttt{Priority\_Resolution.thy}] defines domain-independent locales
  for liveness properties under Byzantine faults: \texttt{priority\_system}
  for deterministic selection from a finite message set with injective
  priorities, \texttt{deadlock\_free\_locking} for timeout-based forced
  release of locks, and \texttt{fair\_leader\_system} for starvation freedom
  under periodic honest leader scheduling.  The locales are parameterized
  over abstract types and impose no dependency on the regulatory domain;
  they are reusable for any distributed system requiring deterministic
  conflict resolution, deadlock freedom, or fair scheduling.

\item[\texttt{DQuencer\_Instance.thy}] instantiates the generic liveness
  framework with the regulatory domain.  It defines authority levels
  (\textsc{Regional}, \textsc{National}, \textsc{International}), a
  four-component priority key with proven injectivity, a Byzantine fault
  model with the BFT threshold $n \geq 3f + 1$, and instantiates each of
  the three liveness locales for the regulatory consensus system.  It
  proves the main liveness theorems
  (\emph{select highest deterministic}, \emph{deadlock freedom},
  \emph{eventual completion}) and concludes with the
  \emph{combined safety liveness} theorem that links the liveness results
  back to the synchronization function defined in
  \texttt{Regulatory\_Instance.thy}, discharging the honest-node assumption
  of the safety proof under Byzantine faults.
\end{description}

\subsection{Theory Dependency Structure}

The four theories form two independent generic-instance pairs that meet
at \texttt{DQuencer\_Instance.thy}:
\texttt{State\_Preservation.thy} is imported by \texttt{Regulatory\_Instance.thy};
\texttt{Priority\_Resolution.thy} depends only on \texttt{Main} and is
deliberately kept independent of the regulatory domain to preserve
reusability; \texttt{DQuencer\_Instance.thy} imports both
\texttt{Priority\_Resolution.thy} and \texttt{Regulatory\_Instance.thy},
serving as the join point where the liveness framework is applied to the
regulatory model and combined with the safety results.

\subsection{Design Decisions}

The regulatory state model incorporates several deliberate design decisions
based on legal precedence and operational considerations:

\begin{itemize}
\item \textsc{Recover} and \textsc{Liquidate} are excluded from the state
  machine as they involve force transfers or external DEX interactions
  rather than state transitions.
\item The direct transition \textsc{Seized} $\to$ \textsc{Frozen} is
  excluded because seizure is a strictly stronger legal constraint than
  freezing; relaxation requires explicit release first.
\item The direct transition \textsc{Frozen} $\to$ \textsc{Restricted} is
  excluded to enforce passage through \textsc{Active}.
\item Preemptive locking prevents concurrent regulatory actions on the
  same asset, \emph{not} double-spending (which is handled at a different
  layer).
\end{itemize}

The liveness model adopts an assume-guarantee approach for starvation
freedom: rather than directly modeling the probabilistic guarantees of
VRF-based leader election (which would require extensive probability
theory infrastructure), the \texttt{fair\_leader\_system} locale assumes
that within any \texttt{fairness\_bound} consecutive epochs, at least one
epoch has an honest leader.  Under $f < n/3$ Byzantine faults, the
probability that a Byzantine leader is elected $k$ consecutive times is
$(f/n)^k$, which decreases exponentially.  The fairness bound captures
this as a deterministic upper bound, following standard assume-guarantee
reasoning practice.

% include generated text of all theories
\input{session}

\section{Acknowledgments}

The regulatory state model (five states, seven actions, and the
transition exclusions) is motivated by the Regulatory Compliance
Protocol (RCP), a framework developed by the Oraclizer research team
that systematically classifies requirements from 15 global financial
regulatory authorities for tokenized capital markets.  RCP is currently
proposed as an Informational Ethereum Improvement Proposal (EIP) in
Draft status.

A companion paper presenting these proofs together with the surrounding
context is available as
\texttt{arXiv:2604.03844}~\cite{kim2026cdsp}, which discusses the
mechanized proofs, the safety-liveness combination structure, the
relation to Lochbihler's Merkle Functor pattern, and the application
context in the regulatory state synchronization domain.

\begin{thebibliography}{1}
\bibitem{kim2026cdsp}
J.~Kim (for the Oraclizer Core Team),
``Safety and Liveness of Cross-Domain State Preservation under Byzantine
Faults: A Mechanized Proof in Isabelle/HOL,''
arXiv:2604.03844 [cs.CR], 2026.
\url{https://arxiv.org/abs/2604.03844}.
\end{thebibliography}

\end{document}
